\documentclass[10pt,a4paper]{article}

% Русский язык и шрифты
\usepackage{fontspec}
\usepackage{polyglossia}
\setdefaultlanguage{russian}

% Шрифт с поддержкой кириллицы
\setmainfont{DejaVu Sans}

% Поля
\usepackage{geometry}
\geometry{
    left=1.5cm,
    right=1.5cm,
    top=1.5cm,
    bottom=1.5cm
}

% Списки
\usepackage{enumitem}

% Заголовки
\usepackage{titlesec}

% Цвета и ссылки
\usepackage{xcolor}
\definecolor{linkbtn}{HTML}{1F5FD6}
\usepackage{hyperref}

\hypersetup{
    colorlinks=true,
    linkcolor=black,
    urlcolor=linkbtn,
    pdftitle={CV - Leonid Nekrylov}
}

\titleformat{\section}
{\large\bfseries}
{}
{0em}
{}
[\titlerule]

\titlespacing{\section}{0pt}{1ex}{1ex}

\setlength{\parindent}{0pt}

% Запись в разделе опыта: название, дата, роль, место одним блоком
\newcommand{\entry}[4]{%
  \par\noindent\textbf{#1}\hfill\textit{#2}\par
  \noindent\textit{#3}\hfill\textit{#4}\par\vspace{2pt}}
\newcommand{\entrygap}{\par\vspace{0.5cm}}

% Ссылка-кнопка: синяя плашка и подчёркнутый адрес рядом
\setlength{\fboxsep}{2.5pt}
\newcommand{\projlink}[3]{%
  \par\noindent\href{#1}{\colorbox{linkbtn}{\textcolor{white}{\small\textbf{▶~#2}}}}%
  \hspace{6pt}\href{#1}{\small\underline{#3}}\par\vspace{3pt}}

\begin{document}

\pagestyle{empty}

% ================= HEADER =================

\begin{center}

{\LARGE \textbf{Леонид Некрылов}}

\vspace{0.1cm}

{\large Robot Learning Engineer}

\vspace{0.2cm}

Москва |
\href{mailto:nekrylov.ln@phystech.edu}{nekrylov.ln@phystech.edu}

\\

\href{https://t.me/leonid_n_n}{t.me/leonid\_n\_n}

\end{center}


% ================= EDUCATION =================

\section*{Образование}

\textbf{Московский физико-технический институт (МФТИ)} \hfill \textit{2023 -- 2027} \\
\textit{Бакалавриат, Физтех-школа прикладной математики и информатики (ФПМИ)} \hfill Долгопрудный, РФ \\
Направление: «Прикладная математика и физика»

\subsection*{Релевантные курсы}

\vspace{0.1cm}
\begin{itemize}[leftmargin=*,nosep]
\item \textbf{Оптимизация и ML:} Годовой курс по методам оптимизации в ML (Д. Меркулов), Машинное обучение-1 (Р. Нейчев), LLM (Giraffe AI)
\item \textbf{Reinforcement Learning:} Самостоятельное освоение курса UC Berkeley CS285 (\textit{Deep Reinforcement Learning}, Sergey Levine).
\item \textbf{Математика:} Теория вероятностей, Случайные процессы, Математическая статистика, Линейная алгебра, Математический и функциональный анализ.
\item \textbf{Computer Science:} Алгоритмы и структуры данных, C++, Базы данных (PostgreSQL), Python.
\item \textbf{Физика/Инженерия:} Введение в робототехнику (О. Буличев), Теоретическая механика.
\end{itemize}


% ================= EXPERIENCE =================

\section*{Опыт работы и проекты}

\entry{\parbox[t]{0.65\textwidth}{UNITAS: исследование 3D World-Action Model для роботических манипуляций}}{Ноябрь 2026 -- настоящее время}{Исследовательский проект}{}
\projlink{https://github.com/leonidnn/unitas-wam-evaluation.git}{Смотреть код}{github.com/leonidnn/unitas-wam-evaluation}
\begin{itemize}[leftmargin=*,nosep]
\item Развернул UNITAS~--- модель на 1,7 млрд параметров для генерации действий робота и прогнозирования движения объектов в 3D.
\item Подготовил воспроизводимую оценку policy в LIBERO: сравнение при разном освещении на одинаковых начальных состояниях, сопоставление checkpoints и измерение задержки инференса.
\item Реализовал запись эпизодов симуляции и локальное воспроизведение в MuJoCo с произвольного ракурса без повторного запуска нейросети.
\item Проверил перенос модели на собственную сцену: подготовил вход из RGB-видео, в котором LLM оценила глубину. Сравнил проекции предсказанных траекторий с отложенными кадрами.
\end{itemize}

\entrygap

\entry{Pet or Threat, поведение четвероногого робота Unitree Go1}{Октябрь 2026}{Пет-проект, MuJoCo}{Москва, РФ}
\projlink{https://leonidnn.github.io/pet-or-threat-site/}{Смотреть видео}{leonidnn.github.io/pet-or-threat-site}
\begin{itemize}[leftmargin=*,nosep]
\item Реализовал автономный выбор поведения четвероногого робота в симуляции без внешних команд: охрана, игровое и социальное взаимодействие, ожидание владельца и отдых.
\item Модуль принятия решений получает данные восприятия и формирует высокоуровневые команды перемещения и направления взгляда; локомоцию обеспечивают готовые RL-политики. Модуль не зависит от симулятора; адаптация к другому роботу выполняется заменой нижнего уровня управления.
\end{itemize}

\entrygap

\entry{ЦКМ МФТИ, RL-агент для схвата объектов по RGB-D}{Сентябрь 2026 -- Октябрь 2026}{Исследовательский проект, Isaac Sim, UR10e}{Москва, РФ}
\projlink{https://github.com/leonidnn/isaac-rgbd-grasping}{Смотреть видео и код}{github.com/leonidnn/isaac-rgbd-grasping}
\begin{itemize}[leftmargin=*,nosep]
\item Построил в Isaac Sim среду для обучения схвату с рандомизацией объектов. Разработал собственный Vulkan-слой для явного выбора GPU без использования остальных устройств на общем многопользовательском сервере.
\item Обучил RL-агента: сеть по карте глубины оценивает успешность схвата для каждой точки и угла, агент выбирает вариант с максимальной оценкой.
\item Проверка в новых случайных позах с полным подводом манипулятора дала 18/20 и 19/20 успешных схватов на двух объектах. SAC не освоил задачу за 5.7 тыс. попыток.
\item Проверил перенос навыка на сложный деформируемый объект (пакет) и сравнил стратегии дообучения с сохранением качества на прежних объектах. Лучший из проверенных вариантов~--- две сети с переключением по типу объекта: 88\% на новом объекте против 42\% без дообучения и 95\% на прежних.
\item Выявил артефакт рендеринга, исказивший обучение, и исправил постановку.
\end{itemize}

\entrygap

\entry{НБ-Технологии, Hybrid Solver}{Сентябрь 2025 -- Ноябрь 2025}{LLM-разработчик}{Москва, РФ}
\projlink{https://huggingface.co/spaces/leonid-nn/hybrid-solver}{Попробовать решатель}{huggingface.co/spaces/leonid-nn/hybrid-solver}
\projlink{https://github.com/leonidnn/hybrid-solver}{Смотреть код}{github.com/leonidnn/hybrid-solver}
\begin{itemize}[leftmargin=*,nosep]
\item Разработал решатель математических задач с проверкой ответа и отказом от выдачи результата, если проверка не пройдена.
\item Использовал принцип перехода от алгоритма Монте-Карло к алгоритму Лас-Вегас (аналогия с классом ZPP): ответ выдаётся только при согласованности результатов трёх независимых методов~--- рассуждения языковой модели, точного вычисления в SymPy и численной проверки.
\end{itemize}

\entrygap

\entry{Проект Cool Deal носимых систем климат-контроля}{Декабрь 2024 -- Сентябрь 2025}{Основатель, тимлид проекта}{Москва, РФ}
\begin{itemize}[leftmargin=*,nosep]
\item Руководил командой из 5 человек, объединяющей разработку и продажи, от идеи до рабочего прототипа носимой системы климат-контроля.
\item Проектировал и собирал прототипы: CAD, датчики, прошивка микроконтроллеров, стендовые испытания.
\item Привлёк грант 1 млн р. «Студенческий стартап» на НИОКР и испытания.
\item Занимался продажами и провёл 2 пилотных внедрения. По результатам проверки гипотезы принял решение о закрытии проекта.
\end{itemize}

\entrygap

\entry{Сбер (Команда рекомендательных систем / Amazme)}{Июль 2024 -- Ноябрь 2024}{Data Scientist}{Москва, РФ}
\begin{itemize}[leftmargin=*,nosep]
\item Собирал витрины на PySpark из сырого кликстрима.
\item Дообучал и валидировал модели отбора кандидатов для рекомендательной системы на PyTorch по офлайн-метрикам.
\item Локализовал причины аномалий в работе рекомендательной системы по сообщениям об ошибках; разрабатывал юнит-тесты.
\end{itemize}

\vspace{0.2cm}


% ================= SKILLS =================

\section*{Технические навыки}

\begin{itemize}[leftmargin=*,nosep]
\item Python (NumPy, SciPy, PyTorch, Matplotlib, PySpark, SymPy), PostgreSQL, Git, Pytest, Linux, \LaTeX, Docker, ROS2 NVIDIA Isaac Sim / Isaac Lab, MuJoCo
\item Английский~-- C1 (TOEFL iBT 107).
\end{itemize}



\end{document}
